\documentclass[12pt]{scrartcl}
\usepackage[sexy]{james}
\usepackage[noend]{algpseudocode}
\setlength {\marginparwidth}{2cm}
\usepackage{answers}
\usepackage{array}
\usepackage{tikz}
\newenvironment{allintypewriter}{\ttfamily}{\par}
\usepackage{listings}
\usepackage{xcolor}
\usetikzlibrary{arrows.meta}
\usepackage{color}
\usepackage{mathtools}
\newcommand{\U}{\mathcal{U}}
\newcommand{\E}{\mathbb{E}}
\usetikzlibrary{arrows}
\Newassociation{hint}{hintitem}{all-hints}
\renewcommand{\solutionextension}{out}
\renewenvironment{hintitem}[1]{\item[\bfseries #1.]}{}
\renewcommand{\O}{\mathcal{O}}
\declaretheorem[style=thmbluebox,name={Chinese Remainder Theorem}]{CRT}
\renewcommand{\theCRT}{\Alph{CRT}}
\setlength\parindent{0pt}
\usepackage{sansmath}
\usepackage{pgfplots}

\usetikzlibrary{automata}
\usetikzlibrary{positioning}  %                 ...positioning nodes
\usetikzlibrary{arrows}       %                 ...customizing arrows
\newcommand{\eqdef}{=\vcentcolon}
\newcommand{\tr}{{\rm tr\ }}
\newcommand{\im}{{\rm Im\ }}
\newcommand{\spann}{{\rm span\ }}
\newcommand{\Col}{{\rm Col\ }}
\newcommand{\Row}{{\rm Row\ }}
\newcommand{\dint}{\displaystyle\int}
\newcommand{\dt}{\ {\rm d }t}
\newcommand{\PP}{\mathbb{P}}
\newcommand{\horizontal}{\par\noindent\rule{\textwidth}{0.4pt}}
\usepackage[top=3cm,left=3cm,right=3cm,bottom=3cm]{geometry}
\newcommand{\mref}[3][red]{\hypersetup{linkcolor=#1}\cref{#2}{#3}\hypersetup{linkcolor=blue}}%<<<changed

\tikzset{node distance=4.5cm, % Minimum distance between two nodes. Change if necessary.
         every state/.style={ % Sets the properties for each state
           semithick,
           fill=cyan!40},
         initial text={},     % No label on start arrow
         double distance=4pt, % Adjust appearance of accept states
         every edge/.style={  % Sets the properties for each transition
         draw,
           ->,>=stealth',     % Makes edges directed with bold arrowheads
           auto,
           semithick}}


% Start of document.
\newcommand{\sep}{\hspace*{.5em}}

\pgfplotsset{compat=1.18}
\begin{document}
\title{CMSC451: Algorithms}
\author{James Zhang\thanks{Email: \mailto{jzhang72@terpmail.umd.edu}}}
\date{\today}

\definecolor{dkgreen}{rgb}{0,0.6,0}
\definecolor{gray}{rgb}{0.5,0.5,0.5}
\definecolor{mauve}{rgb}{0.58,0,0.82}

\lstset{frame=tb,
  language=Java,
  aboveskip=3mm,
  belowskip=3mm,
  showstringspaces=false,
  columns=flexible,
  basicstyle={\small\ttfamily},
  numbers=left,
  numberstyle=\tiny\color{gray},
  keywordstyle=\color{blue},
  commentstyle=\color{dkgreen},
  stringstyle=\color{mauve},
  breaklines=true,
  breakatwhitespace=true,
  tabsize=3
}


\maketitle
These are my notes for UMD's CMSC451: Algorithms. These notes are taken live in class 
(``live-\TeX``-ed). This course is taught by Professor Andrew Childs. 
\tableofcontents

\newpage

\begin{definition}[Algorithm]
  An \vocab{algorithm} is a well-defined computational procedure for processing some input 
  to obtain a desired output.
\end{definition}

\section{Introduction}

\subsection{Stable Matching Problem}

\begin{definition}[Stable Matching Problem]
  The goal is to pair memebers of two sets according to their preferences, such as 
  students and employees for summer internships. Given two sets $S, J$ a \vocab{matching} 
  is a set of pairs of the form $(s,t). s \in S, t \in J$, such that every element of $S$ 
  appears in at most one pair, and similarly for $J$. 
\end{definition}

\begin{note}
For simplicity, the two sets have the same size, and we must match every member of the first set 
with exactly one member of the second set. 
\end{note}

\begin{definition}
  A matching is \vocab{perfect} if every element of $S$ and $J$ is in some pair. 
\end{definition}

\begin{definition}
  A matching is \vocab{stable} if there is no pair $(s,t)$ not in the matching such that 
  $s$ prefers $t$ over its assigned partner and $t$ prefers $s$ over its assigned partner. 
\end{definition}

\begin{example}
  See August 29. notes for some examples. 
\end{example}

\begin{definition}
  Gale-Shapley Algorithm

\begin{lstlisting}
All companies and students are unengaged
While there is a company c that is unengaged and has not made an offer to every students
  Let s be the highest-ranked student for c to whom c has not yet made an offer
  If s is unengaged
    (c, s) become engaged
  Else
    Let c` be the company to which s is engaged
    if s prefers c to c` then 
      (c,s) become engaged
      c` becomes unengaged
    endif
  endif
endwhile
Return the set of engaged pairs
\end{lstlisting}
\end{definition}

\begin{lemma}
  Once a student receives an offer, they remain engaged for the rest of the algorithm, and the 
  quality of their engagement (according to their preference) can only increase. 
  \begin{proof}
    Omitted.
  \end{proof}
\end{lemma}

\begin{lemma}
  As the algorithm proceeds, the sequence of students engaged to a given company can only decrease 
  in performance. 
  \begin{proof}
    Omitted. 
  \end{proof}
\end{lemma}

\begin{lemma}
  If there are $n$ companies and $n$ students, then the algorithm terminates after at most 
  $n^2$ iterations of the while loop. 
  \begin{proof}
    Omitted. 
  \end{proof}
\end{lemma}

\begin{lemma}
  The Gale-Shapley algorithm outputs a perfect matching. 
  \begin{proof}
    Omitted. 
  \end{proof}
\end{lemma}

\begin{theorem}
  The matching output by the G-S algorithm is stable. 
  \begin{proof}
    Omitted.
  \end{proof}
\end{theorem}

\subsection{Asymptotic Notation}

We're interested in asymptotic upper bounds on the worst-case running times of algorithms, 
and we express this using \vocab{Big-O Notation}

\begin{definition}
  $f(n) = \mathcal{O}(g(n))$ if $\exists c > 0, n_0 \geq 0$ such that for all $n \geq n_0$, 
  $f(n) \leq cg(n)$
\end{definition}

\begin{definition}
  $f(n) \in \Omega(g(n))$ if $\exists \ c > 0, n_0 \geq 0$ such that for all $n \geq n_0$, $f(n) \geq cg(n)$. 
\end{definition}

\begin{example}
  $10n^2 \in \mathcal{O}(n^2), \mathcal{O}(n^3), \Omega(n^2), \Omega(n), \cdots$
\end{example}

\begin{definition}
  We say $f(n) \in \Theta(g(n))$ if $f(n) \in \mathcal{O(g(n))}$ and $f(n) \in \Omega(g(n))$
\end{definition}

\begin{definition}
  We say $f(n) \in o(g(n))$ if $\underset{n\to\infty}{\lim}\frac{f(n)}{g(n)} = 0$ and 
  $f(n) \in \omega(g(n))$ if $\underset{n\to\infty}{\lim}\frac{g(n)}{f(n)} = 0$
\end{definition}

\subsection{Basic Graph Algorithms}

\begin{definition}
  A \vocab{directed (undirected)} graph $G = (V, E)$ is a finite set $V$ of vertices and 
  $E$ of edges, where edges are \vocab{ordered (unordered)} pairs of vertices. 
\end{definition}

\begin{definition}
  Vertices are sometimes called \vocab{nodes}
\end{definition}

\begin{definition}
  Edge $\{u, v\}$ joins $u$ and $v$, which are its \vocab{ends}. $u$ and $v$ are 
  \vocab{adjacent}; the edge is \vocab{incident} in $u$ and $v$. Edge $\{u, v\}$ has 
  tail $u$ and head $v$
\end{definition}

\begin{definition}
  The \vocab{degree} of vertex $v$ is the number of incident edges, in a digraph, 
  indeg(v) is the number of edges with $v$ as head, and vice versa for outdeg(v). 
\end{definition}

\begin{note}
  We often write $|V| = n$ and $|E| = m$. Furthermore, note 
  \begin{align*}
    \sum_{v \in V} \deg(v) = 2m && \sum_{v \in V} indeg(V) = m = \sum_{v \in V} outdeg(v) && m \leq \binom{n}{2} = \mathcal{O}(n^2)
  \end{align*}
  in an unidrected graph. In a directed graph, $m \leq n^2$. 
\end{note}

\begin{definition}
  A \vocab{subgraph} of $G = (V, E)$ is a graph $G' = (V', E')$ with $V' \subseteq V$ and 
  $E' \subseteq E$ 
\end{definition}

\subsubsection{Graph Representations}

\begin{definition}
  An \vocab{adjacency matrix} is a $|V| \times |V|$ matrix indexed by vertices, with $A_{uv} = 1$ if $(u, v) \in E$ and $0$ otherwise. 
\end{definition}

\begin{definition}
  A \vocab{adjacency list}: for each $v \in V$, given a list $Adj_v$ of its neighbors outgoing
\end{definition}

\subsubsection{Paths, Cycles, and Connectivity}

\begin{definition}
  A \vocab{path} is a sequence of vertices $(v_0, \cdots, v_k)$ such that 
  $(v_i, v_{i+1}) \in E$ for all $i \in \{0, 1, \cdots, k-1\}$. 
  A path is \vocab{simple} if all of its vertices and edges are distinct. The length is the its number 
  of edges.
\end{definition}

\begin{definition}
  A \vocab{cycle} is a path with $v_0 = v_k$ and all other vertices are distinct. 
\end{definition}

\begin{definition}
  The \vocab{distance} between $u, v \in V$ is the length of the shortest path from $u$ to $v$. 
\end{definition}

\begin{definition}
  An undirected graph is \vocab{connected} if for all $u, v \in V, \ \exists$ a path from $u$ to $v$. A 
  \vocab{component} of a graph is a set of vertices such that pairs of vertices in the set 
  have a path betwen them, and no superset has this property. 
\end{definition}

\begin{definition}
  A digraph is \vocab{strongly connected} if $forall \ u, v \in V, \ \exists$ a path from $u$ to $v$ and a 
  path from $v$ to $u$.
\end{definition}

\begin{definition}
  A \vocab{strongly connected component} is a maximal subgraph that is strongly connected. 
\end{definition}

\subsubsection{Trees}

\begin{definition}
  A \vocab{forest} is a graph with no cycles, and a \vocab{tree} is a connected forest. Note that 
  a \vocab{n-vertex} tree has $n-1$ edges. If we designate one vertex of the tree as a root, then the 
  unique neighbor of a vertex that is closer to the root is the parent, and the neighbors that 
  are further from the root are its chilcren. 

  Vertices of degree $1$ in a tree are called \vocab{leaves}
\end{definition}

\begin{definition}
  A digraph with no cycles is called a \vocab{Directed Acylic Graph} or \vocab{DAG}
\end{definition}

\subsubsection{More on Connectivity}

Given a graph, we want to find is it connected? Given verticecs $s, t$ is there a path 
between them $(s-t)$ connectivity?

One approach is to find a \vocab{spanning tree}, a subgraph that includes all vertices 
and is a tree. 

\begin{proof}[Solution]
  We can specify a rooted tree by giving the parent \verb|pr[v]| of every vertex with 
  \verb|pr[root] = |$\emptyset$. Here is the pseudocode for generating a spanning tree for a tree rooted at $s$. 
  \begin{lstlisting}
    Let T be the graph with one vertes, s, with pr[s] = emptyset. 
    While exists an egde {u, v} joining a vertex u of T to a vertex v not in T
      Add vertex v and edge {u, v} to T
      Set pr[v] = u
  \end{lstlisting}
\end{proof}

The resulting graph is always a tree. 

\begin{lemma}
   The algorithm produces a spanning tree for the component of $s$. 

   \begin{proof}
    Any vertex in the tree has a path to $s$: repeatedly take parents until your reach $s$. 
    To find a path from $u$ to $v$ in $T$, join the paths from $u$ to $s$ and $v$ to $s$. 
    On the other hand, if $u$ is not in $T$, then there is no path to $s$. If there were, we could 
    follow the path and find an edge from a vertex not in $T$ to a vertex in $T$, but no such 
    edge remains when the algorithm is done. 
   \end{proof}
\end{lemma}

\begin{definition}
  \vocab{Breadth-First Search (BFS)}: explore the graph in order of increasing distance from the root.
  \begin{lstlisting}
    Let T be the tree with one vertes, s, with pr[s] = empty
    Repeat:
      Let u be the unexhausted vertex that joined the tree earliest, the "active vertex"
      Choose an edge {u, v} where v is not in T
      Add vertex v and edge {u, v} to T
      Set pr[v] = u
    Until all vertices are exhausted
  \end{lstlisting}
  The total running time of BFS is $\mathcal{O}(\sum_{u \in V} (\deg(u) + 1)) = \mathcal{O}(m + n)$
\end{definition}

\begin{lemma}
  BFS partitions the vertices into layers $L_j$. A vertex in $L_j$ has distance $j$ from $s$. 
  \begin{proof}
    By induction on $j$. Clearly, $L_0 = \{s\}$ is the set of vertices at distance $0$ from $s$.
    If $L_0, L_1, \cdots, L_j$ satisfies the lemma, then vertices in $L_{j+1}$ must have distance at leaast $j+1$ from $s$. 
    But there exists a path of length $j+1$ from any vertex in $L_{j+1}$ to $s$ since there is an edge to a vertex in $L_j$, which has a path of length 
    $j$ to $s$ by IJ. Hence, vertices in $L_{j+1}$ satisfies the lemma, so the lemma follows by induction. 
  \end{proof}
\end{lemma}

\begin{definition}
  A graph $G = (V, E)$ is \vocab{bipartite} if there is a partition $(A, B)$ of $V$ where
  $V = A \cup B, A \cap B = \emptyset$, such that for all $\{u, v\} \in E$, either $u \in A$ and $v \in B$ 
  or $u \in B$ and $v \in A$. 
\end{definition}

\begin{lemma}
  Let $T$ be a BFS tree of graph $G$, and let $u, v$ be vertices of $T$ with $u \in L_i$ and 
  $v \in L_j$. Suppose $\{u, v\}$ is an edge of $G$. Then $|i - j| \leq 1$. 
  \begin{proof}
    WLOG, suppose $u$ is added before $v$. Then $i \leq j$ and $u$ is active before $v$. There are two cases: 
    \begin{enumerate}
      \item $v$ is in the tree when $u$ becomes active. Then $\{u, v\}$ is a non-tree edge, and 
      \verb|pr[v]| joined the tree before $u$, so $j = layer(pr(v)) + 1 \leq i + 1$. 
      \item $v$ is not in the tree when $u$ becomes active. Then $\{u, v\}$ is a tree edge and $j = i + 1$. 
    \end{enumerate}
  \end{proof}
\end{lemma}

\begin{theorem}
  A connected graph with BFS tree $T$ is bipartite if and only if there is no non-tree edge joining vertices in the same layer. 
  \begin{proof}
    Suppose there is no non-tree edge joining vertices in the same layer. 
    Let $A = L_0 \cup L_2 \cup \cdots \cup L_{2k}$ and let $B = L_1 \cup L_3 \cup \cdots \cup L_{2k+1}$. 
    This is a partition since all non-tree edges join vertices whose layers differ by $1$. 

    Conversely, suppose $\{u, v\}$ is a non-tree edge between vertices in the same layer. Suppose their nearest common ancestor 
    is $j$ layers higher, then the path in the tree betwee $u$ and $v$ has length $2j$. Adding $\{u, v\}$ 
    gives a cycle of length $2j + 1$, but an odd cycle cannot be bipartite. 
  \end{proof}
\end{theorem}

\subsubsection{Connectivity in Digraphs}

We can use BFS to find the vertices that can be reached from a given starting vertex. The idea 
is to run BFS fro mevery vertex, so there are $n$ runs of $BFS$. 

\begin{note}
  The digraph is \vocab{strongly connected} if and only if all runs reach all vertices, and the running time would be 
  $\mathcal{O}((m + n)n)$
\end{note}

\begin{lemma}
  If $u$ and $v$ are mutually reachable, and $v$ and $w$ are mutually reachable, then
  $u$ and $w$ are mutually reachable. 
  \begin{proof}
    We can go from $u \to w$ by going from $u \to v \to w$ and we can go from $w \to u$ by going from 
    $w \to v \to u$. 
  \end{proof} 
\end{lemma}

The better idea is the following algorithm
\begin{lstlisting}
  Run a BFS starting from s
  Run BFS starting from s, reversing the directions of all edges
  If both searches find all vertices, the graph is strongly connected, otherwise not
\end{lstlisting}
The running time of this is $\mathcal{O}(2(m + n)) = \mathcal{O}(m + n)$

\subsubsection{Topological Sorting}

The idea here is that we can use a digraph to capture a set of operations with consraints on which ones must be done 
before others. 

\begin{definition}
  A \vocab{topological ordering} is a possible sequence in which the actions can be performed, subject to the contraints on which other actions must have been 
  completed. 
\end{definition}

\begin{lemma}
  If $G$ has a topological ordering, then $G$ is a DAG.
  \begin{proof}
    For a contradiction, suppose $G$ has a topological ordering and a directed cycle $C$, 
    $v_1, v_2, \cdots, v_n$. Let $v_i$ be the lowest-index vertex on $C$, and let $v_j$ be the 
    vertex before it in $C$. Then $(v_j, v_i)$ is an edge with $i < j$, but we also must have $j < i$ for this to be a topological ordering. This 
    is a contradiction. 
  \end{proof}
\end{lemma}

\begin{lemma}
  Every DAG has a vertex with indegree $0$. 
  \begin{proof}
    We proceed by proving the contrapositive: if every vertex has positive indegree, then the 
    graph is a directed cycle. Start at any vertex. Walk along some edge in the backward direction, which must be possible because every vertex has positive indegree.
    Repeat. After at most $n+1$ steps, we must have visited some vertex twice, so we found a cycle. 
  \end{proof}
\end{lemma}

\begin{lemma}
  Every DAG has a topological ordering. 
  \begin{proof}
    By induction on $n = |V|$. Base case, a 1-vertex graph has a topological ordering. Suppose the lemma holds for 
    any digraph with at most $n$ vertices. For an $(n+1)$-vertex graph, let $v$ be a vertex with indegree $0$. 
    Put it first in the ordering. Follow it by a topological oprdering $G \slash \{v\}$. This exists 
    because $G \slash \{v\}$ . is obtained by deleting a vertex and edges from a DAG, and deletion cannot 
    guarantee a cycle, and plus the IH, and we have shown a valid topogical ordering. 
  \end{proof}
\end{lemma}

\begin{theorem}
  Using the above Lemmas, we have shown that a graph is a DAG if and only if it has a 
  topological ordering. 
\begin{lstlisting}
TopoSort(G):
  Let S be an emptyset
  For all vertices v
    If v has indegree 0, add v to S
    Set count[v] = indeg(v)
  endfor
  while S is nonempty do
    remove a vertex v from S
    output v
    for each vertex u that v points to 
      remove edge (v, u) from the graph
      decrement count[u]
      if count[u] = 0, then add u to s
    endfor
  endwhile
\end{lstlisting}
There is $\mathcal{O}(n)$ initialization, the while loop goes through all vertices; for loop for $v$ takes 
$\mathcal{O}(outdeg(v) + 1)$ steps $\implies \mathcal{O}(n + \sum_v (outdeg(v) + 1)) = \mathcal{O}(m + n)$
\end{theorem}

\section{Algorithm Design Techniques}

\subsection{Greedy Algorithms}

A common problem for greedy algorithms is the coin change problem. Fix the coin denominations
$c_1, c_2, \ldots, c_k$, and given a value $v$, find a way to make up $v$ using as few coins as possible. 
\begin{lstlisting}
Change(v):
  If v = 0, return 0
  Else
    Let c be the largest coin of denomination at most v
    Add c to the list of coins
    Change(v-c)
  endif
\end{lstlisting}
The question that we want to ask is, "does this find the maximum number of coins?"

\begin{note}
  If the coin denominations are $1, 2$ then yes. A counterexample $1, 7, 10$ and $v = 14$, as the 
  greedy algorithm would say $10 + 1 + 1 + 1 + 1 = 14 \implies 5$ but the optimal is actually 
  $7 + 7 = 14 \implies 2$
\end{note}

\subsubsection{Shortest Paths}

Given weights $w_{uv} \geq $ for each edge $(u, v) \in E$, what is the length of a shortest path from 
$s$ to $t$, for $s, t \in V$?

\begin{definition}
  \vocab{Dijkstra's Algorithm} finds the shortest path in a digraph with no negative edge weights. 
\begin{lstlisting}
  Let d[s] = 0 and pr[s] = empty
  For all v in V\{s}, ket d[v] = infinity
  Mark all vertices as unexplored
  Repeat:
    Let u be the unexplored vertex with smallest d[u]
    For each neighbor v of u
      If d[u] + w_{uv} < d[v] then 
        Let d[v] = d[u] + w_{uv}
        Let pr[v] = u
      endif
    endfor
    Mark u as explored
  Until all verticse are explored
\end{lstlisting}
\end{definition}

\begin{lemma}
  When $v$ is explored, $d[v]$ is the distance from $s$ to $v$. 
  \begin{proof}
    By induction on the number of explored vertices, base case, this clearly holds when $s$ is explored. 
    IH: suppose this lemma holds when $k$ vertices are explored. Let $v$ be the $(k+1)$st explored vertex; let 
    \verb|pr[v] = u|. By IH, the algorithm has already found the shortest path from $s \to u$. 
    Syppose for a contradiction that the shortest path from $s$ to $v$ does not go through $u$. Let 
    $x$ be the last explored vertex on the path and let $y$ be its unexplored neighbor on this path. Then 
    \[dist(s, v) < d[v] \leq d[y] = \text{length of path from s to y found in alg} \leq dist(s, v)\]
    which is a contradiction.  
  \end{proof}
\end{lemma}

\subsubsection{Minimum Spanning Trees}

Goal: given a graph $GG$ with edge weights, find a spanning tree $T$ that minimizes the total 
weight of edges in $T$

\begin{definition}
  A \vocab{cut} in a graph is a bipartition of the vertices $(V = A \cup B, A \cap B = \emptyset)$. 
  A cut is \vocab{nontrivial} if neither $A$ nor $B$ is empty. An edge \vocab{crosses} the cut 
  if t has one vertex on each side.
\end{definition}

\begin{definition}
  \vocab{Kruskal's Algorithm} adds the lowest-cost edge that doesn't create a cycle in order 
  to construct a MST. 
\end{definition}

\begin{lemma}
  Assume all edge weights are distinct. Let $(A, B)$ be a nontrivial cut; let $e$ be the minimum 
  weight crossing this cut. Then any minimum spanning tree must contain $e$. 
  \begin{proof}
    Let $T'$ be a MST that does not contain $\{u, v\}$. Since $T'$ is spanning, there exists a unique path 
    in $T'$ from $u$ to $v$. Since $u$ and $v$ are on opposite sides of the cut, then there must be an edge along this path
    that crosses the cut; call it $e'$. Let $T$ be the graph obtained from $T'$ by deleting $T'$ and adding $e$. We 
    can show that $T$ is a spanning tree. Since $e$ is the min-weight edge crossing the cut, $w_e < w_{e'}$, 
    and since all other edges are unchanged, $T$ has lower total weight than $T'$. 
  \end{proof}
\end{lemma}
  
\begin{theorem}
  Kruskal's algorithm outputs a MST.
  \begin{proof}
    Suppose the algorithm adds edge $e = \{u, v\}$ to the forest $F$. Consider the cut with $A$ = vertices in the component of $u$ in $F$. 
    Vertex $v$ is on the other side of the cut since adding $e$ doesn't create a cycle. There are no edges in $F$ 
    crossing the cut, so $e$ is in fact the min-weigtht edge crossing the cut, so $e$ is in fact the 
    min-weight edge crossing the cut. So, by the lemma, every MST contains $e$. 
  \end{proof}
\end{theorem}

\begin{definition}
  The idea of \vocab{Prim's Algorithm} is to choose any root and grow a tree by repeatedly adding the non-tree
  vertex with the lowest attachment cost. 
\end{definition}

\begin{theorem}
  Prim's Algorithm outputs a MST.
  \begin{proof}
    Suppose the algorithm adds $e = \{u, v\}$ in the forest $F$. Consider the cut with $A = $ the component of the 
    root in $F$. Clearly, $e$ crosses the cut, and it has the lowest weight of any such edge by the greedy rule, so by the 
    lemma, every MST contains $e$. The algorithm outputs a spanning tree since as long as $F$ is not connected, 
    we can add some edge and still have no cycles. 
  \end{proof}
\end{theorem}

\begin{note}
  Both of these algorithms can be implemented in $\mathcal{O}(m \log n)$ with appropriate 
  data structures. 
\end{note}

\subsubsection{Interval Scheduling}
In this course, we cover three greedy problems with respect to interval scheduling
\begin{enumerate}
  \item finding the largest possible compatible subsets
  \item minimizing the resources used
  \item minimizing lateness
\end{enumerate}

Given $n$ requests with starting times $s_i$ and finishing times $f_i$ for $i \in \{1, \ldots, n\}$. 
A subset of requests is compatibile if no two overlap in time. Find the largest possible subset.

\begin{note}
  When we think about how we want to be greedy, we want to \textbf{choose the earliest-finishing activity first}!
\begin{lstlisting}
GreedyIntervalSchedule(s, f):
  Sort tasks by increasing order of finish time
  Let A be an empty set
  Let f_prev = -inf
  For i = 1 to n:
    if s_i > f_prev 
      Add task i to A
      f_prev = f_i
    endif
  endfor
\end{lstlisting}
The running time of this algorithm is $\mathcal{O}(n \log n)$. Note that the scheduling contains 
no conflicts because we only add intervals that start after the last finished interval. 
\end{note}

\begin{theorem}
  GreedyIntervalSchedule outputs an optimal solution. 
  \begin{proof}
    Let $G = (g_1, \ldots, g_k)$ be the greedy schedule (with indices of intervals in order of increasing finishing time). 
    Let $B = (b_1, \ldots, b_l)$ be an optimal schedule. Let $j$ be the first index where the schedules differ. Thus, 
    \begin{align*}
      G = (g_1, \ldots, g_{j-1}, g_j, \ldots, g_k) && B = (g_1, \ldots, g_{j-1}, b_j, \ldots, b_l)
    \end{align*} 
    By the greedy rule, $f_{g_j} \leq f_{b_j}$. So, let's switch $B$ to the valid schedule
    \[B' = (g_1, \ldots, g_{j-1}, b_{j+1}, \ldots, b_l)\]
    Repeating, we can get another valid schedule 
    \[B'' = (g_1, \ldots, g_k, b_{k+1}, \ldots, b_k)\]
    Since the greedy algorithm stopped with only $k$ jobs, there cannot be a job $b_{k+1}$ that it could have 
    scheduled, so in fact $k = l$ and the greedy schedule is optimal. 
  \end{proof}
\end{theorem}

\subsubsection{Interval Partitioning}

2. Now suppose we need to schedule all the intervals and we want to minimize the resources used. 

\begin{definition}
  The \vocab{depth} of a set of intervals is 
  \[\max_t |\{i \in \{1, \ldots, n\} : t \in [s_i, f_i) \}|\]
  Clearly the number of resources is at least the depth. 

\begin{lstlisting}
GreedyIntervalPartition(s, f):
  Sort the intervals by starting time
  For i = 1 to n:
    Let E = emptyset
    For j = 1 to i - 1:
      If [sj, fj] overlaps [si, fi] then add color[j] to E
    endfor
    Let color[i] be the smallest postiive integer that is not in E
  endfor
\end{lstlisting} 
\end{definition}

\begin{theorem}
  Using GreedyIntervalPartition, each interval is assigned a color from 
  
  $\{1, 2, \ldots, depth\}$ and no two overlapping intervals are the same color. 
  \begin{proof}
    Whereever we assign a color, it is different from the colors of all overlapping intervals, so we never 
    assign the same color to two overlapping intervals. Now we need to show that we don't exclude all 
    colors from $\{1, 2, \ldots, depth\}$. Say that when considering interval $i$, it overlaps $t$ previous intervals. 
    Then $\ \exists \ t + 1$ overlapping intervals, so depth $\geq t + 1 \implies t \leq \text{depth} - 1$, so therer 
    is always a color available. . 
  \end{proof}
\end{theorem}

\subsubsection{Scheduling to Minimize Lateness}

3. Now consider the problem where we are scheduling to minimize lateness. Each task has a duration 
$t_i$ and a deadline $d_i$. Define the \vocab{lateness} of task $i$, if scheduled in 
$[s_i, f_i] = [s_i, s_i + t_i]$ as $l_i = \max\{0, f_i - d_i\}$ and our goal is to minimize
\[\max_i l_i\]
the worst-case lateness. 

\begin{note}
  Here, the strategy is to schedule in order of increasing deadline $d_i$. 
\end{note}

\begin{definition}
   A schedule has an \vocab{inversion} if job $i$ comes before job $j$ but $d_i > d_j$. 
\end{definition}

\begin{lemma}
  All schedules with no inversions and no idle time have the same max lateness.
  \begin{proof}
    The only flexibility in such a schedule is in the order of jobs with the same deadline. Any such 
    jobs must be scheduled consecutively. The finishing time of the last of these jobs doesn't depend on the order. 
  \end{proof}
\end{lemma}

\begin{lemma}
  There is an optimal schedule with no inversions and no idle time. 
  \begin{proof}
    
  \end{proof}
\end{lemma}

\subsection{Divide and Conquer}

\begin{note}
  First, recall mergesort: to sort a list of length $n$, we must divide the list into two halves and sort 
  each of them recursively, and then merge the two lists in time $\mathcal{O}(n)$. The recurrence relation can be modeled as 
  \[T(n) = 2T(\frac{n}{2}) + \mathcal{O}(n) \implies T(n) = \mathcal{O}(n \log n)\]
  by the Master's Theorem
\end{note}

\subsubsection{Closest Points}

Given $n$ points in the plane, find the closest pair. A simple approach finds such a pair in 
$\mathcal{O}(n^2)$ time. 

The main idea of the better algorithm is to recursively divide into left and right halves, find the closest pair 
of points on each side, check whether there is a closer pair of poihnts with one on the left and one on the right. 
As we recursive, maintain sorted lists by x and y coordinates. 

Let us have input points $P = \{p_1, \ldots, p_n\}$ where each $p_i = (x_i, y_i)$. 
\begin{align*}
  Q = \text{pts in P w first n/2 coords} && R = \text{pts in P with last n/2 coods}
\end{align*}
Find the closest pair of points in $Q = (q_0^*, q_1^*)$ and $R = (r_0^*, r_1^*)$. Let 
$\delta = \min(dist(q_0^*, q_1^*), dist(r_0^*, r_1^*))$ and let $x^*$ be the largest x coordinate of a point in 
$Q$. Let $L$ be the line $x = x^*$. 

\begin{lemma}
  If $\exists \ q \in Q, r \in R$ with $dist(q, r) < \delta$, then both $q, r$ must be within distance $\delta$ of $L$. 
\end{lemma}

\begin{lemma}
  If two points in $S$ have distance less than $\delta$ then they must be within $15$ positions of each other in the 
  list of points sorted by $y$ coordinate. 
  \begin{proof}
    Divide the strip into sequences of side length $\frac{\delta}{2}$. 
  \end{proof}
\end{lemma}

\subsubsection{Polynomial Multiplication and the Fast Fourier Transform}

Problem: given polynomials $a(x) = a_0 + a_1x + a_2x^2 + \cdots a_{n-1}x^{n-1}$ and 
$b(x) = b_0 + b_1x + b_2x^2 + \cdots b_{n-1}x^{n-1}$ compute their product 
\[c(x) = \]



\end{document}

